Dans un volume $V_{A}=\qty{50.0}{\mL}$ d'une solution d'acide chlorhydrique de
concentration~:
\[
    C_{A}=\qty{0.010}{\mol\per\L},
\]
on ajoute progressivement, avec une burette graduée, une solution d'hydroxyde de
sodium (ou soude) de concentration~:
\[
    C_{B}=\qty{0.020}{\mol\per\L}.
\]
\begin{enumerate}
    \item Quelle est la réaction de ce dosage~?
    \item Quand dit-on qu'on est à l'équivalence~?
    \item On a déterminé le volume à l'équivalence~:
          \[
              v_{B\mathrm{E}}=\qty{25.0}{\mL}.
          \]
          Ce résultat est-il conforme aux valeurs attendues~?
\end{enumerate}

